\begin{block}{Mecanismo Seesaw Quiral}
  La supresión natural de la masa exige la inclusión de un término de Majorana
  para la componente dextrógira estéril. El lagrangiano extendido es
  \begin{equation}
    \mathcal{L}_{\text{Mass}} = \overline{\nu}_{L}M_{D}\nu_{R} 
    + \frac{1}{2}\overline{\nu_{R}^{c}}M_{R}\nu_{R} + h.c.
  \end{equation}

  Bajo la condición de una masa activa intrínseca nula ($M_L \approx 0$) y 
  definiendo la masa de Dirac quiral como $M_D = m e^{-i\alpha}$, la dinámica 
  queda determinada por la matriz de masa
  \begin{equation}
    \mathcal{M} = \begin{pmatrix} 
      0 & m e^{-i\alpha} \\ 
      m e^{-i\alpha} & M_R 
    \end{pmatrix}.
  \end{equation}

  En el régimen de alta energía ($M_R \gg m$), la diagonalización de la matriz 
  revela dos estados físicos con valores propios bien diferenciados:
  \begin{equation}
    m_1 \approx \frac{m^2}{M_R} e^{-2i\alpha}, \quad m_2 \approx M_R.
  \end{equation}

  Aquí, $m_1$ corresponde al neutrino activo levógiro, mientras que $m_2$ 
  representa la excitación dextrógira estéril. 
  El ángulo quiral $\alpha$ determina unívocamente la fase del neutrino activo, 
  constituyendo una fuente geométrica directa de violación CP leptónica.
\end{block}